%%=============================================================================
%% DISCUSION
%%=============================================================================
\chapter{Discusión}
\label{cap:discusion}

Este capítulo interpreta los resultados del Capítulo~\ref{cap:resultados}, los contrasta con la literatura revisada en el Estado del arte y discute sus implicaciones para el uso de QPP-RNG en dispositivos IoT críticos.

\section{Lo que las pruebas estadísticas pueden y no pueden ver}

El hallazgo más importante de este trabajo es metodológico. En dos escenarios distintos, la salida de QPP-RNG superó las pruebas estadísticas estándar mientras su fuente aportaba mucha menos entropía de la que el diseño contabiliza. En los microcontroladores (E1), la fuente no aportó nada: la salida es idéntica tras cada reinicio y, aun así, supera el nivel 1 y obtiene casi 5 bits por muestra con los estimadores de SP 800-90B. En los anfitriones rápidos (E4), la fuente aporta unos 2 bits por byte, pero la salida obtiene entre 7.2 y 7.4 bits por byte en SP 800-90B y, acondicionada, supera todas las pruebas.

La explicación es la misma en ambos casos. QPP-RNG es un sistema con realimentación: la medición de tiempo se convierte en semilla, la semilla determina los \textit{pads} y los \textit{pads} determinan la siguiente medición. Un sistema así puede producir una secuencia de aspecto aleatorio aunque la entrada externa sea nula o escasa, del mismo modo que un generador pseudoaleatorio produce una secuencia de aspecto aleatorio a partir de una semilla fija. Las pruebas de SP 800-22, ENT e incluso los estimadores de SP 800-90B aplicados a la salida miden cómo se ve la secuencia, no cuánta información impredecible entró en ella. Esto es coherente con lo que el propio SP 800-90B establece: la estimación de entropía debe hacerse sobre la fuente de ruido, antes de cualquier procesamiento determinístico \cite{Turan2018RecommendationGeneration}.

Estos resultados matizan los de la evaluación publicada de QPP-RNG \cite{Vrana2025RawEvaluation}, que reporta valores de min-entropía cercanos al máximo sobre la secuencia de salida en múltiples plataformas. La reproducción de este trabajo coincide con esas cifras cuando se mide lo mismo: la salida cruda de la referencia obtiene entre 7.2 y 7.4 bits por byte. La diferencia aparece al medir otra cosa, la entropía fresca de $\Delta t$ condicionada a $n_p$. No es posible afirmar que la implementación original de los autores se comporte exactamente igual, porque su código no está publicado. Sin embargo, la propiedad que explica el déficit es estructural y no depende de detalles de implementación: el tiempo de una caminata de $n_p$ pasos es, en primera aproximación, proporcional a $n_p$, y $n_p$ está determinado por la semilla. Cualquier implementación que mida su propio bucle hereda esa relación, y solamente el residuo respecto de ella es entrada nueva.

La consecuencia práctica es una regla de diseño: en un generador basado en \textit{jitter}, la entropía debe estimarse en la fuente, con el estado previo como condición, y la evaluación de la salida debe tratarse como una verificación necesaria, pero nunca suficiente. La comparación entre reinicios desde el mismo estado inicial es una prueba barata y de alto valor que las baterías estándar no incluyen.

\section{Por qué la referencia falla en microcontroladores}

La literatura atribuye la entropía de QPP-RNG a fenómenos como la caché, la predicción de saltos, la planificación del sistema operativo y la contención por recursos \cite{Kuang2025IID-BasedJitter,Vrana2025RawEvaluation}. Ninguno de ellos existe en un AVR o en un Cortex-M4 que ejecuta desde Flash sin sistema operativo: cada instrucción tarda un número fijo de ciclos, y el temporizador cuenta el mismo reloj que la CPU. Por lo tanto, el resultado de E1 no es una anomalía de las placas evaluadas, sino la consecuencia esperable del modelo de ejecución de esta clase de procesadores, que es precisamente la que domina en el IoT restringido.

Este modo de falla es especialmente peligroso porque es silencioso. Las pruebas de salud continuas de SP 800-90B no lo detectan, pues la fuente varía de una muestra a otra; las pruebas estadísticas tampoco. Un dispositivo desplegado con la referencia generaría, tras cada reinicio, la misma secuencia de claves, que cualquiera que conozca el \textit{firmware} podría reproducir. Dada la importancia de la unicidad del \textit{nonce} en los esquemas AEAD (Sección~\ref{sec:aead}), esto equivale a reutilizar claves y \textit{nonces} entre dispositivos y entre reinicios.

La solución que se propone no consiste en más cómputo, sino en un reloj físicamente independiente. El perro guardián de los ATmega, alimentado por su propio oscilador RC, ofrece entre 1.6 y 2.2 bits por muestra sin hardware adicional, y el RNG del nRF52840 ofrece 7.24 bits por byte. Ambos cambian la naturaleza del problema: la entropía proviene de un fenómeno físico, el ruido de fase de un oscilador o el ruido térmico, y no de la complejidad de la microarquitectura.

\section{El papel de la estructura de QPP en la variante IoT}

En la variante IoT, la caminata de QPP deja de ser la fuente de entropía y pasa a ser una etapa determinística de mezcla. Los resultados lo confirman: en el nRF52840, la min-entropía estimada pasa de 7.24 bits por byte en el RNG crudo a 7.28 en la salida, es decir, la caminata no añade entropía medible ni la destruye. Su efecto observable es estadístico: los bytes crudos del RNG fallan la prueba de rachas y la salida la supera. En los AVR ocurre lo mismo con una fuente mucho más pobre: las muestras del perro guardián, con 1.6 a 2.2 bits cada una y una deriva lenta que las hace fallar gravemente las pruebas de nivel 1, se convierten en una salida que supera la batería completa sobre 1\,000\,000 de bytes con 7.26 a 7.36 bits por byte en \texttt{ea\_non\_iid}. Esa mejora no proviene de la caminata, sino de que cada byte de salida reúne 8 muestras, es decir, al menos 12.9 bits evaluados; la caminata solo distribuye esa entropía de manera uniforme sobre los 8 bits del byte.

Esto obliga a una conclusión honesta. Desde el punto de vista de la entropía, una tubería formada solamente por la fuente física, las pruebas de salud y el acondicionador SHA-256 ofrecería las mismas garantías, a menor costo. En el nRF52840, la caminata limita el rendimiento a 101 B/s, mientras que el periférico RNG por sí solo entrega del orden de un kilobyte por segundo. En AVR, en cambio, la caminata optimizada termina antes de que llegue la siguiente muestra del perro guardián, de modo que su costo no reduce el rendimiento. La estructura de QPP se conservó porque es el objeto de estudio de esta tesis y porque, con $N = 6$ y cuatro caminatas por byte, es barata en las plataformas donde la fuente es lenta. Pero su conservación es una decisión de diseño, no un requisito de seguridad, y debería reconsiderarse en plataformas con una fuente rápida.

\section{Contabilidad de entropía y márgenes}

La variante IoT hace explícito lo que la referencia deja implícito: cuántos bits de min-entropía se acreditan a cada muestra y, por lo tanto, cuántas muestras se necesitan por byte. Ese número controla tanto la tasa de salida como los umbrales de las pruebas de salud, y su correcta elección es, en la práctica, la principal decisión de seguridad de quien integra la librería. El experimento E2 mostró por qué los márgenes importan: las estimaciones con un millón de muestras fueron entre 24 y 25\,\% menores que con 20\,000, porque la deriva lenta del oscilador solo se hace visible en capturas largas. Un valor acreditado a partir de capturas cortas habría quedado demasiado cerca de la estimación final. El valor de 1 bit por muestra en AVR conserva un margen de $1.6\times$ en el peor de los dos ejemplares evaluados, pero ese margen corresponde a un ejemplar y a unas condiciones de temperatura y tensión; un despliegue real debe reevaluar la fuente en el modelo de dispositivo concreto.

\section{Rendimiento suficiente para sembrar, no para todo}

Con 3.7 B/s de salida acondicionada, la tubería en AVR no está pensada para producir un \textit{nonce} aleatorio por mensaje. Su uso adecuado es sembrar, al arranque, un generador determinístico criptográficamente seguro, por ejemplo un DRBG de SP 800-90A \cite{Barker2015RecommendationGenerators} o una construcción basada en la función de salida extensible de Ascon \cite{Turan2025Ascon-BasedFunctions}, del que luego se derivan claves y \textit{nonces}. Una semilla de 256 bits toma unos 9 segundos, un costo aceptable en el arranque de un sensor, pero no en cada operación. Para los \textit{nonces} de Ascon-AEAD128 en particular, un contador persistente que sobreviva a los cortes de energía ofrece una garantía de unicidad que ningún generador aleatorio de 128 bits puede dar por sí solo.

\section{Rust para fuentes de entropía}

Rust cumplió su papel en la mayor parte de la librería: el mismo código \texttt{no\_std} compila para el anfitrión, para Cortex-M y para AVR; los genéricos con parámetros constantes ($N$, el generador interno y el temporizador) no tienen costo en tiempo de ejecución; y el sistema de tipos permitió expresar como un valor de retorno (\texttt{HealthError}) una falla de la fuente que en otras implementaciones pasa inadvertida. Sin embargo, el trabajo también mostró un conflicto entre un compilador optimizador y una fuente de entropía que depende del tiempo de su propio código: la optimización completa colapsó el generador de referencia en un punto fijo. La mitigación (fijar el nivel de optimización de dos \textit{crates} y usar barreras \texttt{black\_box}) funciona, pero depende de una configuración que el usuario de la librería podría alterar sin advertirlo. La variante IoT elimina este conflicto de raíz, porque su entropía no depende del tiempo de ejecución de su propio código.

\section{Lecciones sobre el proceso de validación}

Dos episodios del trabajo muestran que la validación de una fuente de entropía también puede fallar por causas externas al generador. El primero fue la falla de la prueba de independencia por $\chi^2$ en las muestras cortas de AVR: sin los controles de \texttt{/dev/urandom} del mismo tamaño, habría sido fácil atribuirla al generador. El segundo fue la instalación de STS mediante un enlace simbólico, que produjo informes de SP 800-22 sin sentido pero de apariencia plausible. Ambos se detectaron porque se aplicó la misma regla: comparar cada herramienta y cada tamaño de muestra contra un caso de resultado conocido antes de interpretar el caso nuevo. En una evaluación de seguridad, la cadena de herramientas forma parte de lo que hay que validar.
